\section{תנע זוויתי, סיבובים וקוונטיזציה בשדה אלקטרומגנטי}

סיכום זה מכסה שני נושאים מרכזיים: השלמת הדיון הפורמלי באלגברת התנע הזוויתי (כולל מטריצות סיבוב ופונקציות ויגנר), ומעבר לטיפול בחלקיק טעון בשדה אלקטרומגנטי באמצעות הפורמליזם הלגרנז'יאני וההמילטוניאני לקראת קוונטיזציה.

\subsection{אלגברת התנע הזוויתי והצגות בלתי פריקות}

בשיעורים הקודמים פיתחנו את הקשר בין חבורת הסיבובים \LR{$SO(3)$} לבין האלגברה שלה. ראינו שכל סיבוב $R$ ניתן לכתיבה באמצעות האקספוננט של היוצרים (\LR{generators}) של החבורה:
\begin{equation}
    R(\vec{\theta}) = \exp\left(-\frac{i}{\hbar} \vec{\theta} \cdot \vec{J}\right)
\end{equation}
כאשר האופרטורים $J_i$ מקיימים את יחסי החילוף של האלגברה:
\begin{equation}
    [J_i, J_j] = i\hbar \epsilon_{ijk} J_k
\end{equation}
(ביחידות שבהן $\hbar=1$, כפי שהופיעו בהרצאה, הפקטור הוא $i$).

\subsubsection{סולם המצבים ופעולת אופרטורי הסולם}
הגדרנו את האופרטורים המועילים:
\begin{itemize}
    \item $J^2 = J_x^2 + J_y^2 + J_z^2$ (אופרטור קזימיר, חילופי עם כל יוצרי האלגברה).
    \item $J_\pm = J_x \pm i J_y$ (אופרטורי סולם - העלאה והורדה).
\end{itemize}

הבסיס הסטנדרטי למרחב ההילברט (או לתת-המרחב האינווריאנטי) מוגדר על ידי הוקטורים העצמיים המשותפים של $J^2$ ו-$J_z$:
\begin{equation}
    J^2 |j,m\rangle = j(j+1) |j,m\rangle, \quad J_z |j,m\rangle = m |j,m\rangle
\end{equation}

כדי לחשב אלמנטי מטריצה של סיבובים כלליים, עלינו לדעת כיצד פועלים אופרטורי הסולם על המצבים. נשתמש בזהות הבאה:
\begin{equation}
    J_- J_+ = J^2 - J_z^2 - \hbar J_z
\end{equation}
(ובאופן דומה עבור $J_+ J_-$ עם סימן הפוך באיבר האחרון).

נניח פעולה מהצורה $J_+ |j,m\rangle = c_+ |j,m+1\rangle$. מתוך חישוב הנורמה $\langle \psi | \psi \rangle$ עבור $|\psi\rangle = J_+ |j,m\rangle$:
\begin{equation}
    |c_+|^2 = \langle j,m | J_- J_+ | j,m \rangle = \langle j,m | (J^2 - J_z^2 - J_z) | j,m \rangle
\end{equation}
התוצאה היא $j(j+1) - m(m+1)$.

\begin{definitionbox}{בחירת הפאזה (\LR{Condon-Shortley Phase})}
נהוג לבחור את המקדמים $c_\pm$ כרכיבים \textbf{ממשיים וחיוביים}. בחירה זו נקראת \LR{Condon-Shortley convention}. תחת בחירה זו, פעולת אופרטורי הסולם היא:
\begin{align}
    J_+ |j,m\rangle &= \sqrt{j(j+1) - m(m+1)} |j,m+1\rangle \\
    J_- |j,m\rangle &= \sqrt{j(j+1) - m(m-1)} |j,m-1\rangle
\end{align}
\end{definitionbox}

\subsubsection{סיבובים סופיים: זוויות אוילר ופונקציות $D$ של ויגנר}
סיבוב כללי במרחב התלת-ממדי מתואר על ידי שלוש זוויות אוילר $(\alpha, \beta, \gamma)$. הקונבנציה המקובלת במכניקת הקוונטים (סיבוב $z-y-z$) מגדירה את אופרטור הסיבוב כך:
\begin{equation}
    D(\alpha, \beta, \gamma) = e^{-i J_z \alpha / \hbar} e^{-i J_y \beta / \hbar} e^{-i J_z \gamma / \hbar}
\end{equation}
משמעות הסיבוב (בגישה האקטיבית על המערכת):
\begin{enumerate}
    \item סיבוב סביב ציר $z$ בזווית $\gamma$.
    \item סיבוב סביב ציר $y$ (החדש) בזווית $\beta$.
    \item סיבוב סביב ציר $z$ (החדש) בזווית $\alpha$.
\end{enumerate}



אלמנטי המטריצה של אופרטור הסיבוב בבסיס $|j,m\rangle$ נקראים \textbf{מטריצות ויגנר} (\LR{Wigner D-matrices}):
\begin{equation}
    D^j_{m'm}(\alpha, \beta, \gamma) = \langle j,m' | e^{-i J_z \alpha} e^{-i J_y \beta} e^{-i J_z \gamma} | j,m \rangle
\end{equation}
מכיוון שהמצבים הם עצמיים ל-$J_z$, ניתן להוציא את התלות ב-$\alpha$ ו-$\gamma$ החוצה כפאזות:

\begin{resultbox}{פונקציות ויגנר}
\begin{equation}
    D^j_{m'm}(\alpha, \beta, \gamma) = e^{-i(m'\alpha + m\gamma)} d^j_{m'm}(\beta)
\end{equation}
כאשר $d^j_{m'm}(\beta)$ היא "מטריצת ויגנר הקטנה" המממשת סיבוב סביב ציר $y$:
\begin{equation}
    d^j_{m'm}(\beta) = \langle j,m' | e^{-i J_y \beta} | j,m \rangle
\end{equation}
\end{resultbox}

הפונקציות $d^j_{m'm}(\beta)$ ניתנות לחישוב באמצעות טור סופי (נוסחת ויגנר), והן קשורות לפולינומי יעקובי ולפונקציות הרמוניות ספריות. בפרט, עבור המקרה $j=l$ שלם ו-$m=0$, מקבלים קשר ישיר להרמוניות הספריות $Y_{lm}$.

\subsection{דינמיקה של חלקיק בשדה אלקטרומגנטי}

כדי לבצע קוונטיזציה נכונה למערכת הכוללת שדות אלקטרומגנטיים, עלינו להתחיל מהתיאור הקלאסי (לגרנז'יאן) ולעבור בצורה קנונית להמילטוניאן.

\subsubsection{הלגרנז'יאן האלקטרומגנטי}
עבור חלקיק בעל מסה $m$ ומטען $q$ הנע בשדות הנוצרים על ידי פוטנציאל וקטורי $\vec{A}(\vec{x},t)$ ופוטנציאל סקלרי $\phi(\vec{x},t)$, הלגרנז'יאן הוא:
\begin{definitionbox}{לגרנז'יאן של חלקיק בשדה א"מ}
\begin{equation}
    L(\vec{x}, \dot{\vec{x}}, t) = \frac{1}{2}m\dot{\vec{x}}^2 + q \dot{\vec{x}} \cdot \vec{A} - q\phi
\end{equation}
איבר האינטראקציה תלוי במהירות באופן ליניארי.
\end{definitionbox}

\subsubsection{משוואות התנועה (כוח לורנץ)}
נפעיל את משוואת אוילר-לגרנז':
\begin{equation}
    \frac{d}{dt} \left( \frac{\partial L}{\partial \dot{x}_i} \right) - \frac{\partial L}{\partial x_i} = 0
\end{equation}
\begin{enumerate}
    \item התנע הקנוני (נגזרת לפי מהירות):
    \begin{equation}
        \frac{\partial L}{\partial \dot{\vec{x}}} = m\dot{\vec{x}} + q\vec{A}
    \end{equation}
    \item הנגזרת השלמה בזמן:
    \begin{equation}
        \frac{d}{dt}(m\dot{x}_i + qA_i) = m\ddot{x}_i + q\left(\frac{\partial A_i}{\partial t} + (\vec{v}\cdot\nabla)A_i\right)
    \end{equation}
    \item נגזרת לפי מקום (גרדיאנט):
    \begin{equation}
        \frac{\partial L}{\partial x_i} = q \frac{\partial}{\partial x_i} (\vec{v}\cdot\vec{A}) - q \frac{\partial \phi}{\partial x_i}
    \end{equation}
\end{enumerate}
בעזרת זהות וקטורית $\nabla(\vec{v}\cdot\vec{A}) = (\vec{v}\cdot\nabla)\vec{A} + \vec{v}\times(\nabla\times\vec{A})$ (כאשר $\vec{v}$ משתנה בלתי תלוי במקום בגזירה זו), ושימוש בהגדרות השדות:
\begin{equation}
    \vec{E} = -\nabla\phi - \frac{\partial \vec{A}}{\partial t}, \quad \vec{B} = \nabla \times \vec{A}
\end{equation}
מקבלים את כוח לורנץ:
\begin{resultbox}{כוח לורנץ}
\begin{equation}
    m\ddot{\vec{x}} = q(\vec{E} + \vec{v} \times \vec{B})
\end{equation}
זה מוכיח שהלגרנז'יאן הנ"ל אכן מתאר את הדינמיקה הנכונה.
\end{resultbox}

\subsubsection{המעבר להמילטוניאן וקוונטיזציה}
ההמילטוניאן מוגדר על ידי התמרת לז'נדר $H = \vec{p}\cdot\dot{\vec{x}} - L$.
חשוב להבחין בין שני סוגי תנע:
\begin{notebox}{תנע קנוני מול תנע קינטי}
\begin{itemize}
    \item \textbf{תנע קינטי (מכני):} $\vec{\pi} = m\vec{v}$. זהו הגודל הפיזיקלי המדיד (למשל $mv$).
    \item \textbf{תנע קנוני:} $\vec{p} = \frac{\partial L}{\partial \dot{\vec{x}}} = m\vec{v} + q\vec{A}$. זהו הגודל הצמוד לקואורדינטה בפורמליזם ההמילטוני.
\end{itemize}
\end{notebox}

נבודד את המהירות: $\vec{v} = \frac{1}{m}(\vec{p} - q\vec{A})$. נציב בהגדרת ההמילטוניאן:
\begin{equation}
    H = \frac{1}{2m}(\vec{p} - q\vec{A})^2 + q\phi
\end{equation}

\subsubsection{קוונטיזציה קנונית}
כלל הקוונטיזציה הבסיסי קובע שעלינו להחליף את התנע ה\textbf{קנוני} $\vec{p}$ (ולא את הקינטי!) באופרטור הדיפרנציאלי:
\begin{equation}
    \vec{p} \to -i\hbar \nabla
\end{equation}
לכן, משוואת שרדינגר לחלקיק בשדה אלקטרומגנטי היא:
\begin{resultbox}{משוואת שרדינגר בשדה א"מ}
\begin{equation}
    i\hbar \frac{\partial \psi}{\partial t} = \left[ \frac{1}{2m} \left( -i\hbar\nabla - q\vec{A} \right)^2 + q\phi \right] \psi
\end{equation}
\end{resultbox}

\textbf{נקודה חשובה:} בניגוד למכניקה הקלאסית, שבה הפוטנציאלים $\vec{A}, \phi$ הם כלי עזר מתמטי והגדלים הפיזיקליים הם השדות $\vec{E}, \vec{B}$, במכניקת הקוונטים הפוטנציאלים מופיעים ישירות במשוואת התנועה (בהמילטוניאן). הדבר מוביל לאפקטים קוונטיים ייחודיים כגון אפקט אהרונוב-בוהם (\LR{Aharonov-Bohm Effect}), שבהם הפוטנציאל משפיע גם באזורים שבהם השדות מתאפסים, דרך הפאזה של פונקציית הגל.